% !TEX root = ../main.tex
\chapter{はじめに}\label{chap:intro}
% ===== 下書き(Claude、11/20)。レオが読んで理解 → 確認の問い → 直す =====
% TODO(レオ):「なぜファンの稼働を減らすか」(1.3 節)と、題目が決まったら直す

\section{研究背景}
IoT や 5G の普及により、データの発生源の近くで処理を行うエッジコンピューティングの重要性が高まっている\cite{shi2016edge}。エッジでの処理の基盤として、Raspberry Pi などのシングルボードコンピュータを複数台束ねたマイクロクラスタが用いられている\cite{iridispi2014}。マイクロクラスタは安価で、台数を増減しやすく、教育・研究やエッジでの推論、センサデータの処理などに使われている。

一方、エッジの機器は、サーバ室のような専用の空調がない場所に置かれることが多い。オフィスや教室、工場、屋外の箱などである。周囲の環境の温度は機器の温度と性能に大きく影響する\cite{edgeengine2023}。Raspberry Pi 4 Model B(以下、Raspi)は、CPU 温度が約 80\,℃ を超えると動作クロックを下げて発熱を抑える(サーマルスロットリング)。スロットリングが起きると処理が遅くなり、エッジでのリアルタイムの処理に影響する。

\section{熱の問題とファンによる冷却}
温度の上昇を抑える最も直接的な方法は、ファンで冷やすことである。Raspi には、PoE+ HAT やケースに付属するファンを取り付けられる。ファンを常に回し続ければ温度は十分に下がるが、ファンは機械的に動く部品であり、回し続けるほど軸受けが摩耗して寿命が縮む。また、人のいる場所に置かれるエッジ機器では、ファンの騒音も問題となる。したがって、スロットリングを避けられる範囲で、ファンはできるだけ回さないことが望ましい。
% 要:ファンの寿命と稼働時間の関係の文献
% TODO(レオ):ファンを減らす理由はこれでよいか(寿命・騒音を主、電力は大きなファンへの応用)

OS 標準のファン制御は、温度が一定の値を超えたら回す、しきい値による制御である。Raspi の Ubuntu の標準設定では、約 40\,℃ から段階的にファンが回り始める。これは安全側の設定であり、80\,℃ よりはるかに低い温度でファンが回る。必要なときだけファンを回すには、温度がこの先どうなるかを予測し、80\,℃ に達しそうなときだけ回せばよい。

\section{空間熱干渉}
複数のノードを近くに並べると、あるノードの排熱が隣のノードを温める。本研究ではこれを空間熱干渉と呼ぶ。データセンタでは、排熱が吸気側に回り込む熱の再循環が知られている\cite{tang2008sensor,moore2005cool}。空間熱干渉があると、自ノードの負荷だけでは温度の変化を説明できず、温度の予測を誤る原因となる。

先行発表\cite{suzuki2026ipsj}では、囲った環境に Raspi を並べ、片方だけに負荷をかけると、負荷をかけていない隣のノードの温度が約 12\,℃ 上がることを観測した。そこで、隣のノードからの熱の影響を状態として推定するカルマンフィルタを提案した。ただし、後の見直しで、評価の方法にいくつかの課題が見つかった(第\ref{chap:preliminary}章)。また、干渉がどのような設置で生じるかは確かめていなかった。
% TODO(Claude):suzuki2026ipsj を references.bib に追加する

\section{研究目的}
本研究の目的は、Raspi のマイクロクラスタで、スロットリングを避けつつファンの稼働をできるだけ減らすファン制御を実現し、実機で評価することである。そのために、隣のノードからの熱の影響を状態として推定する拡張熱モデルとカルマンフィルタで 30 秒先の温度とその不確かさを予測し、不確かさに応じた安全閾値を予測温度が超えるときだけファンを回す制御を提案する。

評価では、次の3つの問いに答える。

\begin{description}
  \item[RQ1] 対象とする環境で、空間熱干渉はどの程度生じるか
  \item[RQ2] 提案手法は、別の日のデータに対して 30 秒先の温度をどの程度の誤差で予測でき、不確かさは実際の誤差を表しているか
  \item[RQ3] 予測に基づくファン制御は、80\,℃ を超えずにファンの稼働をどれだけ減らせるか。従来の方式と比べて差があるか
\end{description}

先行発表からの主な変更点は、(1)予測に未来の実測値を用いない、同定と評価に別のデータを用いるなど評価の方法を直したこと、(2)ファンをモデルの入力に加え、実際にファンを動かして閉ループで評価したこと、(3)干渉が生じない環境での結果も示し、手法が有効な条件を明らかにすることである。

\section{本論文の構成}
本論文の構成は次のとおりである。第\ref{chap:background}章では、本研究に関わる技術として、Raspi の熱特性、ファン、熱 RC モデル、カルマンフィルタについて述べる。第\ref{chap:related}章では関連研究を述べ、本研究の位置づけを示す。第\ref{chap:preliminary}章では、先行発表の実験と課題、および囲いのない環境とラックでの熱干渉の確認実験について述べる。第\ref{chap:method}章では提案手法を、第\ref{chap:impl}章ではその実装を述べる。第\ref{chap:eval}章で評価実験の結果を示し、第\ref{chap:discussion}章で考察する。最後に第\ref{chap:conclusion}章で本論文をまとめる。
